\newpage
\phantomsection
\label{prf:closure-under-formation-propositional-formulas}
\LRAProofFor{lem:closure-under-formation-propositional-formulas}

\begin{remark*}[Return]
\hyperref[lem:closure-under-formation-propositional-formulas]{Return to Theorem}
\end{remark*}

\ProofVaultURL{https://github.com/wsollers/lra-proof-vault/tree/master/volume-i/chapter-01-propositional-logic/lem-closure-under-formation-propositional-formulas}

\begin{lemma*}[Closure Under Formation]
The set \(\WFF\) is closed under the propositional formation rules.

That is, if \(\varphi\in\WFF\), then
\[
\neg\varphi\in\WFF.
\]
If \(\varphi,\psi\in\WFF\) and
\[
\circ\in\{\land,\lor,\to,\leftrightarrow\},
\]
then
\[
(\varphi\circ\psi)\in\WFF.
\]
\end{lemma*}

\begin{proof}[Professional Standard Proof]
\LRAProofBodyStart
This is exactly the closure part of the recursive definition of
\(\WFF\). If \(\varphi\in\WFF\), then the negation formation rule gives
\(\neg\varphi\in\WFF\). If \(\varphi,\psi\in\WFF\) and
\(\circ\in\{\land,\lor,\to,\leftrightarrow\}\), then the binary formation
rule for \(\circ\) gives \((\varphi\circ\psi)\in\WFF\). Hence \(\WFF\) is
closed under the propositional formation rules.
\end{proof}

\begin{proof}[Detailed Learning Proof]
\LRAProofBodyStart
We prove the two closure clauses separately.

First let \(\varphi\in\WFF\). By the definition of well-formed formulas, the
negation formation rule says that whenever a formula belongs to \(\WFF\), the
string obtained by prefixing it with \(\neg\) also belongs to \(\WFF\). Applying
that rule to \(\varphi\), we obtain
\[
\neg\varphi\in\WFF.
\]

Second let \(\varphi,\psi\in\WFF\), and let
\[
\circ\in\{\land,\lor,\to,\leftrightarrow\}.
\]
The binary formation rule says that whenever two formulas belong to \(\WFF\),
joining them with any allowed binary connective produces another well-formed
formula. Applying that rule to the formulas \(\varphi,\psi\) and to the chosen
connective \(\circ\), we obtain
\[
(\varphi\circ\psi)\in\WFF.
\]

Thus \(\WFF\) is closed under negation and under each of the binary formation
operations \(\land,\lor,\to,\leftrightarrow\).
\end{proof}

\begin{remark*}[Proof structure]
The proof is a direct unpacking of the recursive definition of
\(\WFF\): use the unary formation clause for negation, then use the binary
formation clause for each permitted binary connective.
\end{remark*}

\begin{dependencies}
\begin{itemize}
  \item \hyperref[def:well-formed-formula]{Well-Formed Formula}
\end{itemize}
\end{dependencies}

\clearpage
